% Tapis de Quillen, with Quillen's offprint and preprint (cote 162-5, typed part): archive cote 162-5, site page 68 (PDF page 69).
% preprint. Draft compared: none.
% AI-transcribed by Claude Opus 5.5 (claude-opus-5-5, adaptive thinking, effort high), Message Batches API, batch t162_full (msgbatch_01UAh3zyWPd454AeDkhnmQ33), 2026-11-13; not yet checked by eye.
\scanpage{162-5/68}
\mspage{- 12 -}

for suitable $k,m,n > 0$ . Hence $cl(E_2) = n\,cl(L^k) - cl(L^k \unsure{\otimes|\oplus} E_3^{\unsure{*}})$
and $x = cl(E_1) + cl(L^k \unsure{\otimes|\oplus} E_3^{\unsure{*}}) - n(m - cl(E_4)) = cl(E_1 \unsure{\oplus|\otimes} L^k \unsure{\otimes|\oplus} E_3^{\unsure{*}} \unsure{\oplus|\otimes} nE_4)$
$- nm$ .
% NOTE: a smudged typed dashed line runs across the page below this line; no characters can be read in it

\uline{Bibliography}

[1] Adams, J. F., On the groups J(X)-I, Topology Vol. 2 (1963)
p.p. 181-195.

[2] Artin, M. and Mazur, B., Homotopy of Varieties in the
Etale Topology (to appear).

[3] Artin, M. and Grothendieck, A., S\hand{\'e}minare de G\hand{\'e}om\hand{\'e}trie
algebrique (63-64) - Cohomologie \hand{\'e}tale des Schemas.
(Mimeographed notes, Institute des Hautes \hand{\'E}tudes \hand{'}
Scientifiques, 1963-64).

[4] Hall, I. M., The Generalized Whitney Sum, Quart. J. Math.
Oxford (2), 16 (1965), 360-84.

Massachusetts Institute of Technology

\margin{68}
